\section{Functions, sections, and smoothing on the initial manifold}
\label{sec:preliminaries}

Throughout this Section, distances, balls, norms, and volume are those of
the complete initial metric $g$.  Fix $o\in M$ and put $r(x)=\dist_g(o,x)$.
Constants may depend on $g$ and $o$.  In particular, the estimates below
are not estimates uniform over all choices of center.

We prepare the functions and estimates used in the flow and disc
arguments.  Holomorphic interpolation supplies global functions and
tangent fields for controlling analytic discs, and canonical sections
with growth bounds for the flow's potential barriers.  Heat smoothing
produces a smooth exhaustion with controlled complex Hessian; combined
with the canonical sections, its cutoffs yield averaged curvature
estimates.

\begin{proposition}[Initial analytic data]\label{prop:initial-data}
The following objects and estimates are available.
\begin{enumerate}[label=\textup{(\roman*)}]
\item There is a smooth strictly plurisubharmonic function $\psi$ with
      $\psi\le C(1+r)$.
\item Holomorphic functions and holomorphic tangent fields interpolate
      arbitrary finite jets at a finite set of points.  There are finitely
      many holomorphic functions defining an injective immersion
      $F:M\longrightarrow\C^N$, and finitely many global holomorphic
      tangent fields spanning $T^{1,0}M$ at every point.
\item At every point $x$ there is a global holomorphic canonical section
      $S$ with $S(x)\ne0$ and
      $\log|S|_g^2\le C_S(1+r)$.  There is also a smooth function $\Xi$
      such that
      \begin{equation}\label{prelim:xi}
      \dd\Xi\ge\Ric(g),\qquad \Xi\le C(1+r^{3/2}).
      \end{equation}
\item There is a smooth exhaustion $u\ge1$ such that
      \begin{equation}\label{prelim:exhaustion}
      C^{-1}(1+r)\le u\le C(1+r),\qquad
      |\partial u|_g\le C,\qquad |\dd u|_g\le C/u.
      \end{equation}
\item If $S(o)\ne0$ is chosen as in \textup{(iii)}, then, for $L\ge1$,
      \begin{equation}\label{prelim:averages}
      \frac{1}{\Vol B(o,L)}\int_{B(o,L)}
           \bigl|\log|S|_g^2\bigr|\,dV_g\le CL,
      \qquad
      \frac{1}{\Vol B(o,L)}\int_{B(o,L)}R_g\,dV_g\le C/L.
      \end{equation}
\end{enumerate}
Neither $F$ nor $\psi$ is asserted to be proper.  The exhaustion in
\textup{(iv)} has a two-sided complex Hessian bound.
\end{proposition}

\subsection{The weight and holomorphic interpolation}

Let $\gamma$ be a unit-speed ray starting at $o$ and let
\[
 b(x)=\lim_{j\to\infty}\{d_g(o,\gamma(j))-d_g(x,\gamma(j))\}.
\]
The triangle inequality gives monotonicity of the expressions in braces,
their boundedness on compact sets, and $|b(x)-b(y)|\le d_g(x,y)$.
We recall the complex Hessian comparison that makes $b$ strictly
plurisubharmonic; this is the Busemann input in the argument of
Chen--Zhu \cite[Section~2]{ChenZhu}.

Take $x$ in a fixed compact set and a minimizing segment of length $l$
from $x$ to $\gamma(j)$.  For an initial vector $X$, use its parallel
translate multiplied by $1-\sigma/l$ as a comparison variation field,
and use the same field construction for $JX$.  The second variation of
length is computed on their perpendicular components.  Its curvature
contribution is the negative of
\[
 \int_0^l(1-\sigma/l)^2
 \{\Rm(\dot\gamma,X_\parallel,X_\parallel,\dot\gamma)
   +\Rm(\dot\gamma,JX_\parallel,JX_\parallel,\dot\gamma)\}
 \,d\sigma.
\]
Removing radial components does not change this expression.  It is a
positive multiple of the bisectional curvature before the minus sign.
On a fixed short initial subsegment, positivity is bounded below for
all unit initial vectors and all unit initial directions based in the
chosen compact.  The derivative term in the index form is $O(l^{-1})$.
Consequently the sum of the two second derivatives of the distance is
at most $-c_K|X|^2+O(l^{-1})|X|^2$, with $c_K>0$.

This comparison does not require a smooth distance function.  Vary the
initial endpoint along a holomorphic germ and use the comparison paths
to obtain smooth upper supports for distance.  The sum of the endpoint
acceleration terms vanishes for the two real directions of that germ.
Their negatives give lower supports for the signed distance, with the
same strict Levi lower bound.  An upper test for the signed distance
also lies above its lower support and therefore has this Levi lower
bound.  Local uniform passage to $b$ gives the same viscosity inequality.
Equivalently, testing on complex lines gives the subharmonic
inequality.  Thus $b$ is strictly plurisubharmonic, with a positive Levi
lower bound on each compact.  Richberg approximation, with pointwise
error less than one, gives a smooth strictly plurisubharmonic $\psi$
with $\psi\le1+r$ after adding a constant
\cite{Richberg}\cite[Section~4, Corollary~2]{GreeneWu}.

Here are the bundle choices for the interpolation step.  Write
$K=\bigwedge^nT^{*(1,0)}M$.  A function, a canonical section, and a
tangent field can be regarded respectively as a holomorphic top form
with coefficients in
\begin{equation}\label{prelim:coefficient-bundles}
 K^{-1},\qquad \boldone,\qquad T^{1,0}M\otimes K^{-1}.
\end{equation}
The first has nonnegative curvature, the second is flat, and the last
has nonnegative Nakano curvature: it is $E\otimes\det E$ for the
Griffiths-positive bundle $E=T^{1,0}M$, so the
Demailly--Skoda tensor-with-determinant theorem applies
\cite[Theorem~2.6]{Demailly}.  These statements refer to the coefficient bundles
in \eqref{prelim:coefficient-bundles}; no Nakano positivity of $E$
itself is being assumed.

For completeness, prescribe jets of order $m$ at a finite set $A$.
Choose disjoint coordinate balls and local holomorphic sections with
those jets, multiply them by cutoffs equal to one on smaller balls,
and call the resulting smooth section $s_0$.  Its $\dbar$ is supported
in a compact set disjoint from $A$.  Use the weight
\[
 k\psi+\sum_{a\in A}c_m\vartheta_a\log|z_a|^2,
 \qquad c_m>n+m,
\]
where each $\vartheta_a$ is one near $a$ and supported in the chosen
coordinate ball.  The negative Hessian terms of the cutoff logarithms
are supported in a fixed compact.  A sufficiently large $k$ makes the
total weight strictly plurisubharmonic and dominates those terms.
The inverse-curvature norm of $\dbar s_0$ is integrable, since its
support avoids the poles.  The complete K\"ahler $L^2$ estimate for
bundle-valued top forms gives $v$ with
$\dbar v=\dbar s_0$ and finite weighted norm
\cite{Hormander}\cite[Theorem~5.1]{Demailly}.  The singular-weight statement follows by
regularizing the logarithms, using uniform estimates, and taking a
weak limit.  Near a point of $A$, $v$ is holomorphic; integrability
against $|z_a|^{-2c_m}$ forces its Taylor coefficients through order
$m$ to vanish.  Hence $s_0-v$ has the prescribed jets.  This also proves
point separation by including two points in $A$.

For a canonical section, the same construction gives
\[
 \int_M |S|_g^2 e^{-k\psi}\,dV_g<\infty.
\]
The cutoff pole weights can be dropped in this estimate at the price
of a constant.  The function $|S|_g^2$ is subharmonic: away from its
zeros, $\dd\log|S|_g^2=\Ric(g)$, and its squared norm has the same
nonnegative distributional Laplacian across the zeros.  The local
mean-value inequality on a radius-one ball
\cite[Theorem~2.1]{LiSchoen}, together with the upper bound for
$\psi$, yields
\[
 |S(x)|_g^2\le
 \frac{C_S}{\Vol B(x,1)}\exp\bigl(C_S(1+r(x))\bigr).
\]
Nonnegative Ricci curvature gives
$\Vol B(x,1)\ge c(1+r(x))^{-2n}$ by Bishop--Gromov relative volume
comparison with a fixed ball at $o$ \cite[\S2.1]{GromovBetti}.
Its polynomial loss is absorbed by the exponential.
This proves the growth assertion in Proposition~\ref{prop:initial-data};
compare the canonical-section construction in
\cite[Section~3]{ChenZhu}.

We explain why finite tuples suffice, using the dimension-count
argument underlying parametric transversality.
The compact-open spaces of holomorphic functions and sections are
Fr\'echet and hence Baire spaces. Fix a compact family of points,
or of distinct pairs of points. Surjectivity of finite-jet evaluation,
openness of matrix rank, and a finite cover give a finite-dimensional
space of perturbations whose evaluation derivative is onto on a
neighborhood of that compact family. This derivative is independent
of the perturbation parameter, since the evaluations are affine
linear. Thus the universal evaluation is a submersion there.

For each smooth stratum of the forbidden locus, its inverse image
has real codimension equal to that of the stratum. If this codimension
exceeds the real dimension of the point or pair space, the inverse
image has dimension less than the perturbation space. Its projection
to that space has measure zero: cover it by countably many compact
coordinate patches; on each patch the projection is Lipschitz, and
covering a $d$-dimensional cube by $O(\delta^{-d})$ cubes shows that
its image in $\R^p$, $p>d$, has volume at most $C\delta^{p-d}$.
Consequently arbitrarily small perturbations avoid the forbidden
strata on the prescribed compact family.

For $N$ functions, the rank-$r$ stratum of $N\times n$ derivative
matrices has complex codimension $(N-r)(n-r)\ge N-n+1$ when $r<n$.
For $N\ge2n+1$ this exceeds $n$.
The equal-values condition on distinct pairs has complex codimension
$N>2n$. Avoidance on a compact is open in the compact-open topology,
using Cauchy estimates for derivatives. A countable exhaustion of
$M$ and of $M\times M\setminus\operatorname{diag}$, followed by Baire
intersection, therefore gives one finite injective immersion.
For $m$ tangent sections the least deficient-rank codimension is
$m-n+1>n$ when $m\ge2n$. The same zeroth-jet argument gives a finite
spanning family. Neither conclusion asserts properness.

Choose countably many canonical sections $S_i$ with no common zero.
Write $a_i=|S_i|_g^2$.  The bound
$a_i\le\exp(C_i(1+r))$ implies
\[
 a_i\le e^{B_i}\exp(1+r^{3/2})
\]
for a finite $B_i$, because a linear function of $r$ is bounded above
by $r^{3/2}$ plus a constant.  Choose $c_i>0$ so small that
$c_ie^{B_i}\le2^{-i}$ and that every derivative of $c_i a_i$ of order
at most $i$ has norm at most $2^{-i}$ on the first $i$ exhaustion
compacts.  The sum $\sum c_i a_i$ converges smoothly locally, is
positive, and is bounded above by $\exp(1+r^{3/2})$.  In a local
canonical frame it is $\det(g)^{-1}$ times a sum of squares of
holomorphic functions.  Thus
$\Xi=\log\sum c_i|S_i|_g^2$ satisfies \eqref{prelim:xi}.

\subsection{Hodge heat without a curvature upper bound}

To smooth the distance truncations used below, we need to preserve an
upper bound for their complex Hessians.  Scalar heat estimates will
then control the trace and give the missing lower Hessian bound.
We also record how heat evolution of a bounded closed $(1,1)$-form
produces a scalar potential for its change; this identity will be used
again when constructing the flow comparison cutoff.

Let $H_a$ denote scalar heat for $\Delta_g$.  Normalize the Hodge heat
operator $\mathcal H_a$ on real $(1,1)$-forms so that it agrees with
$H_a$ after taking the trace.  The normalization is also chosen so
that its generator on closed real $(1,1)$-forms is
$A\mapsto\dd\tr_g A$.

The scalar heat kernel has total mass one.  Indeed, fix its pole $y$ and
let $\chi_R$ be a Lipschitz distance cutoff equal to one on $B_g(y,R)$,
supported in $B_g(y,2R)$, and satisfying $|\nabla\chi_R|\leq C/R$.
The Gaussian estimate of Li--Yau \cite{LiYau} and the kernel gradient estimate of
Souplet--Zhang \cite[Theorem~1.3]{SoupletZhang}, together with
Bishop--Gromov comparison \cite[\S2.1]{GromovBetti}, give
\[
 \int_{R\leq d_g(x,y)\leq2R}|\nabla_xH_s(x,y)|\,dV_g(x)
 \leq C s^{-1/2}(1+R/\sqrt{s})^N e^{-cR^2/s}
\]
for a fixed dimensional exponent $N$.  Integrating the heat equation
against $\chi_R$ from $\varepsilon$ to $t\leq T$ and then letting
$\varepsilon\downarrow0$ therefore yields
\[
 \left|\int_M\chi_R(x)H_t(x,y)\,dV_g(x)-1\right|
 \leq \frac{C}{R}\int_0^T
 s^{-1/2}(1+R/\sqrt{s})^N e^{-cR^2/s}\,ds.
\]
The right side tends to zero as $R\to\infty$.  Positivity and the
sub-Markov property of the minimal kernel allow passage to the total
mass, proving $H_t1=1$.  The cutoffs here require only completeness and
the first derivative of distance; they precede the smooth exhaustion
constructed below.

\begin{lemma}[Heat order and commutation]\label{lem:hodge-heat}
For bounded real $(1,1)$-forms, Hodge heat is well defined by cutoff
limits of the complete $L^2$ semigroup.  It preserves order and obeys
\[
 -Cg\le A\le Cg\quad\Longrightarrow\quad
 -Cg\le\mathcal H_aA\le Cg,
 \qquad \tr_g\mathcal H_aA=H_a\tr_g A.
\]
If $A$ is also smooth and closed, then
\begin{equation}\label{prelim:hodge-identity}
 \mathcal H_aA
   =A+\dd\int_0^aH_b\tr_g A\,db.
\end{equation}
If a Lipschitz function $f$ is constant off a compact set and
$\dd f\le Cg$ as currents, then
\begin{equation}\label{prelim:hessian-preservation}
 \dd H_af\le Cg\qquad(a>0).
\end{equation}
\end{lemma}

\begin{proof}
First take smooth compactly supported data.  On a relatively compact
smooth domain use the componentwise zero-Dirichlet realization of
the Hodge heat equation.  Its Weitzenb\"ock formula is a connection
Laplacian plus a zeroth-order curvature reaction.  At a null vector
$X$ of a nonnegative Hermitian form $A$, diagonalize $A$ in a unitary
frame.  The reaction on $(X,\bar X)$ is a positive fixed multiple of
\[
 \sum_j R_g(X,\bar X,e_j,\bar e_j)A_{j\bar j}\ge0;
\]
the Ricci-times-$A$ terms vanish there.  The tensor maximum principle
therefore preserves the nonnegative cone.  Applying the same argument
to $Cg-A$, with nonnegative boundary data, proves the upper bound by
$Cg$.  The trace satisfies the scalar heat equation, including the
zero boundary condition.  This is the null-eigenvector mechanism in
\cite[Section~2]{NiTamStructure}.

We specify the complete realization and the exhaustion step.
On compactly supported smooth forms put $D=d+d^*$.
Completeness gives essential self-adjointness of $D^2$ and hence of
$D$: a deficiency vector for $D$ would lie in the kernel of
$(D^2)^*+1$, which is zero
\cite[Corollary~2.10 and Remark~2.12]{BravermanMilatovicShubin}.
The closed quadratic form of the resulting Hodge Laplacian is
\[
 q(A)=\|dA\|_2^2+\|d^*A\|_2^2,
 \qquad
 V=\overline{C^\infty_c}^{\,(\|\cdot\|_2^2+q)^{1/2}}.
\]
For a nested smooth exhaustion $\Omega_j$, let $V_j$ be the same
closure of forms supported in $\Omega_j$, extended by zero.
On a fixed domain, local ellipticity identifies this with the
full zero boundary-value section domain. It is not an absolute
or relative Hodge boundary condition. The spaces $V_j$ are nested
and their union is dense in $V$.
For $\lambda>0$, the domain resolvent applied to an $L^2$ form,
and extended by zero, is the orthogonal projection of the complete
resolvent onto $V_j$ for the inner product $q+\lambda\langle\cdot,\cdot\rangle$:
both satisfy the same variational identity against $V_j$.
The projections therefore converge in that norm. Strong resolvent
convergence and the common lower bound zero imply strong heat
convergence \cite[Section~6.6, Theorem~6.31 and Corollary~6.33]{Teschl2009}.
Here $\Delta_j$ denotes the nonnegative domain Hodge Laplacian.
Zero extension of its heat means $e^{-a\Delta_j}P_j$, where $P_j$
is restriction followed by zero extension; it does not mean
extending its generator by zero.
The cone and trace assertions pass to this limit and then to bounded
data by cutoff and local regularity.  More explicitly, the complete
Hodge semigroup has the off-diagonal estimate
\begin{equation}\label{prelim:off-diagonal}
 \|\boldone_E\mathcal H_a\boldone_F\|_{2\to2}
 \le C\exp\bigl(-c\dist(E,F)^2/a\bigr).
\end{equation}
To see that this estimate needs no curvature upper bound, use the
self-adjoint first-order operator $d+d^*$.  Its principal symbol has
norm equal to the covector norm, so the wave evolution has finite
propagation speed \cite[Theorem~1.1 and Section~4]{McIntoshMorris}.
The Gaussian Fourier representation of
$\exp(-a(d+d^*)^2)$ then gives \eqref{prelim:off-diagonal}, with the
inessential change of constants required by our normalization
\cite[Theorem~3.4]{CoulhonSikora}.
Since $\Vol B(o,R)\le CR^{2n}$, dyadic annular decomposition shows
that the cutoff evolutions of a bounded form are Cauchy in $L^2$ on
each compact, uniformly for bounded positive times.  Interior
parabolic estimates, whose coefficients are bounded on the compact
under consideration, give smooth convergence away from $a=0$.

Closedness is proved using the complete semigroup, not the
Dirichlet operators on domains.  Let $\zeta_R$ be a compactly supported
Lipschitz cutoff equal to one on $B(o,R)$, supported on $B(o,2R)$,
and satisfying $|d\zeta_R|\le C/R$.  The complete Hodge semigroup
commutes with $d$ on the $L^2$ domains, by the self-adjoint de Rham
complex.
The fixed time normalization does not affect the spectral identity
$D e^{-aD^2}=e^{-aD^2}D$, and projection to degree one higher gives
the $d$ identity. The data $\zeta_RA$ used here belong to the full
$D$ graph domain by compactly supported $H^1$ approximation.
For bounded closed $A$,
\[
 d\mathcal H_a(\zeta_RA)
   =\mathcal H_a^{(3)}(d\zeta_R\wedge A).
\]
The data on the right have $L^2$ norm at most $CR^{n-1}$ and support
outside $B(o,R)$.  Estimate~\eqref{prelim:off-diagonal} in degree
three shows that they tend to zero on each fixed compact, uniformly
for $0<a\le a_0$.  Distributional passage to the limit proves
$d\mathcal H_aA=0$.  Local regularity also gives the initial value
$A$ on each compact.  K\"ahler commutation on closed forms now gives
\[
 \partial_a\mathcal H_aA=\dd\tr_g\mathcal H_aA
                         =\dd H_a\tr_g A.
\]
Integration proves \eqref{prelim:hodge-identity}, first in
distributions and then as an identity of smooth forms for $a>0$.

The passage from closed form heat to a scalar potential by integrating
its trace is also the mechanism in
\cite[Lemma~2.1 and proof of Theorem~2.1]{NiTamHodge2013}.

Finally put $f_0=f-c$, where $c$ is the constant value at infinity.
This is a compactly supported Lipschitz function.  Complete $L^2$
K\"ahler commutation, smooth approximation, and duality give, for a
compactly supported nonnegative Hermitian test form $B$,
\[
 \langle\dd H_af_0,B\rangle
       =\langle\dd f_0,\mathcal H_aB\rangle.
\]
The metric pairing here identifies the cone of nonnegative Hermitian
forms with its dual; equivalently one pairs the current with the
Hodge star of $B$.  Choose $0\le\zeta\le1$ compactly supported and
equal to one near $\supp f_0$.  Since $\dd f_0$ is supported there,
\[
 \langle\dd f_0,\mathcal H_aB\rangle
 =\langle\dd f_0,\zeta\mathcal H_aB\rangle
 \le C\int\zeta\tr_g\mathcal H_aB\,dV_g
 \le C\int\tr_gB\,dV_g.
\]
Positivity and scalar trace commutation justify the last two steps.
Finally $H_a1=1$ under nonnegative Ricci curvature, so
$\dd H_af=\dd H_af_0$.  This proves
\eqref{prelim:hessian-preservation}.
\end{proof}

\subsection{Dyadic smoothing and averaged curvature}

For $L=2^j$, $j\ge0$, choose a fixed smooth nondecreasing scalar
truncation of $r^2/L^2$ which is zero on $r\le L$ and one on
$r\ge2L$, and denote it by $f_L$.  It is Lipschitz, with constant
$C/L$.  The paired index-form comparison with linear parallel fields
gives $\dd r^2\le Cg$ away from the cut locus and in the upper-support
sense at the cut locus.  The derivative-square term introduced by the
truncation is bounded on its annulus.  Hence
\[
 \dd f_L\le CL^{-2}g
\]
as currents.  The word ``smooth'' here describes the scalar
truncation; $f_L$ itself need not be smooth at the cut locus.

Let $w_L=H_{\epsilon L^2}f_L$, with a fixed sufficiently small
$\epsilon>0$.  Gaussian heat bounds and the Lipschitz estimate give
\[
 \|w_L-f_L\|_\infty
 \le (C/L)\sup_x\int d_g(x,y)H_{\epsilon L^2}(x,y)\,dV_g(y)
 \le C\sqrt\epsilon.
\]
Nonnegative Ricci curvature is sufficient for these scalar heat
estimates.  Applying the Li--Yau gradient inequality to the positive
heat solutions $1+H_af_L$ and $2-H_af_L$ gives, at
$a=\epsilon L^2$,
\[
 |\Delta_g w_L|\le C_\epsilon L^{-2},\qquad
 |\partial w_L|_g\le C_\epsilon L^{-1}.
\]
Indeed the first translate bounds $\partial_aH_af_L$ from below,
the second bounds it from above, and substituting both time bounds
in the same inequalities bounds the gradient; see
\cite[Theorem~1.3(i)]{LiYau}.
Lemma~\ref{lem:hodge-heat} gives the upper bound for each complex
Hessian eigenvalue, and the lower trace bound then gives the lower
bound for each eigenvalue.  Thus $|\dd w_L|_g\le C_\epsilon L^{-2}$.

Choose a fixed smooth nondecreasing step $\theta$ which is zero near
zero and one near one, with the constant regions wide enough to
contain the above heat error.  Then $v_L=\theta(w_L)$ is identically
zero on $r\le L$ and one on $r\ge2L$, and
\[
 |\partial v_L|_g\le C/L,\qquad |\dd v_L|_g\le C/L^2.
\]
The locally finite sum
\[
 u=1+\sum_{j\ge0}2^jv_{2^j}
\]
is comparable to $1+r$.  At any point, only a bounded number of
summands have nonzero derivatives; their scales are comparable to
$1+r$.  This proves \eqref{prelim:exhaustion}.  In particular, scalar
cutoffs $p(u/L)$ have gradient $O(L^{-1})$ and absolute complex
Hessian $O(L^{-2})$, with support and transition region in fixed
multiples of $B(o,L)$.

Put $w=\log|S|_g^2$ for a canonical section with $S(o)\ne0$.
The Poincar\'e--Lelong formula gives
\begin{equation}\label{prelim:canonical-current}
 \dd w\ge\Ric(g)
\end{equation}
as currents.  The truncations $w_m=\max(w,-m)$ are subharmonic.
The heat submean inequality and the linear upper growth give
\[
 w_m(o)\le H_{L^2}w_m(o),\qquad
 H_{L^2}(w_m)^+(o)\le CL.
\]
One can justify the submean inequality first on exhaustion domains
and then by Gaussian decay; the positive part has at most linear
growth.  For $m$ large, $w_m(o)=w(o)$, so
$H_{L^2}(w_m)^-(o)\le CL+|w(o)|$.  The heat-kernel lower bound on
$B(o,L)$ is $c/\Vol B(o,L)$.  Monotone convergence as $m\to\infty$
therefore proves the first estimate in \eqref{prelim:averages}.

Take a nonnegative cutoff equal to one on $B(o,L)$ and supported in
$B(o,CL)$ with absolute Laplacian at most $C/L^2$, as just constructed.
Integrating the trace of \eqref{prelim:canonical-current} against it
and moving $\Delta_g$ onto the cutoff gives
\[
 \int_{B(o,L)}R_g\,dV_g
 \le \frac{C}{L^2}\int_{B(o,CL)}|w|\,dV_g
 \le \frac{C}{L}\Vol B(o,L).
\]
The last step uses volume doubling.  This proves the second estimate
and completes Proposition~\ref{prop:initial-data}.  The average
scalar-curvature conclusion is the estimate of
\cite[Theorem~2]{ChenZhu}; the details above record the particular
cutoffs and absolute logarithmic estimate needed below.
